\section{Construction of Quantum Isomorphic Graphs from the Penrose Dodecahedron}

\subsection{Kochen-Specker Sets}

\jvinline{
    Picking up from the $E_8$ subsection right above,
    state why Kochen-Specker sets are interesting here
}
\jvinline{Another example (next to $E_8$) of a KS set is the Penrose Dodecahedron}

\subsection{The Penrose Dodecahedron}

Main sources~\cite{Vlasov2024,ZIMBA1993697}.

Let $\omega = \frac{-1 + \sqrt{3}i}{2}$ be the third root of unity over $\mathbb{C}$, satisfying $\omega^3 = 1$.
The Penrose Dodecahedron can then be defined as the following set of 40 vectors in $\mathbb{C}^4$:
\begin{align*}
    \left(1, 0, 0, 0\right) & \left(0, 1, 0, 0\right) & \left(0, 0, 1, 0\right) & \left(0, 0, 0, 1\right) \\
    \left(\right) & \left(\right) & \left(\right) & \left(\right) \\
    \left(\right) & \left(\right) & \left(\right) & \left(\right) \\
\end{align*}



% \jvinline{Give construction of it, its group structure etc.}

\question{
    Is the group we want to analyse the one called $\mathcal{W}$ in~\cite[Section 1]{Vlasov2024},
    i.e. the "transitive group of 25920 unitary projective transformations of $\mathbb{CP}^3$
    sending the configuration of 40 points to itself"?

}

\jvinline{
    Find central type subgroup with coisotropic stabilizers?
    \begin{itemize}
        \item Unlike for $E_8$, the Pauli Matrices are probably not the solution
        \item Rather, the matrices $R_k \in M_4(\mathbb{C})$, $k \in \{1, \dots, 4\}$ appearing in~\cite{Vlasov2024} seem promising:
        They fulfill the following properties: \begin{itemize}
            \item $R_k$ is unitary, i.e. $R_k R_k^\dagger = R_k^\dagger R_k = 1$
            \item $R_k^3 = 1$
            \item $R_k = R_k^T$
            \item Combining these facts, we get $R_k^2 = R_k^\dagger = \overline{R_k}$ (?)
        \end{itemize} 
        \item Furthermore, in~\cite{ZIMBA1993697} they mention two graphs $Q, Q^*$, which are similar but not isomorphic. Maybe good candidates?
    \end{itemize}
}

\jvinline{
    The approach (as for $E_8$): Look at $G$, construct magic unitary $u$, compute $A_H = u A_G u^\dagger$
    to get the pair $(G, H)$ and hope $G \not \cong H$.
}
